\section{结论}
本文以统一的地形、动态载荷、能量与资源占用模型连接四问。Q1用精确集合划分得到\QOneTrips{}架次；Q2多邻域搜索按延误优先选出\QTwoTrips{}架次；Q3外层搜索与条件LP形成\QThreeTrips{}运输、\RelayTrips{}中继的正式候选；Q4固定其任务时序，精确得到两组、三组分别缺少\GapTwo{}、\GapThree{}份资源。延误改善伴随运输架次和能耗增加，不是所有指标同时改善，外层搜索也没有全局最优证书。

独立加密验证限定了结论：Q3在1 s采样下通过，0.5 s却发现\ViolationsHalf{}个未覆盖状态，最低余量\MarginHalf{} dB；最小硬截止余量为零。当前联合候选尚需修复通信缺口，Q4仅为其固定输出下的条件配置分析。后续应优先补足沿程通信验证与覆盖修复，再扩大离散搜索和开展真实能耗校准。
